\subsection{ソフトウェアアーキテクチャ記述}

\subsubsection*{アーキテクチャ記述の役割}
\addcontentsline{toc}{subsubsection}{アーキテクチャ記述の役割}
小さなシステムを一人で作る場合、アーキテクチャは頭の中にあるだけでも作業できることがある。しかし、大規模で複雑なシステムをチームで作り、運用し、保守する場合には、頭の中の理解だけでは不十分である。人が入れ替わり、要求が変わり、構成が進化するため、アーキテクチャを具体的な成果物として残す必要がある。

この成果物を\keyterm{アーキテクチャ記述}\en{architecture description}という。アーキテクチャ記述は、開発者にとっては構築の設計図であり、テスト担当、品質保証担当、認証担当、運用担当、保守担当にとっては、自分の作業を進めるための根拠である。

かつては、アーキテクチャ記述は文章と非形式的な図で行われることが多かった。しかし、利害関係者と関心事が多様になるにつれて、単一の図ですべてを表すことは難しくなった。そのため、目的や読者に応じて複数のビューを使う方法が重要になった。
